% Options for packages loaded elsewhere
\PassOptionsToPackage{unicode}{hyperref}
\PassOptionsToPackage{hyphens}{url}
%
\documentclass[
  12pt,
]{book}
\usepackage{lmodern}
\usepackage{setspace}
\usepackage{amssymb,amsmath}
\usepackage{ifxetex,ifluatex}
\ifnum 0\ifxetex 1\fi\ifluatex 1\fi=0 % if pdftex
  \usepackage[T1]{fontenc}
  \usepackage[utf8]{inputenc}
  \usepackage{textcomp} % provide euro and other symbols
\else % if luatex or xetex
  \usepackage{unicode-math}
  \defaultfontfeatures{Scale=MatchLowercase}
  \defaultfontfeatures[\rmfamily]{Ligatures=TeX,Scale=1}
  \setmainfont[]{Nimbus Roman}
  \setsansfont[]{Nimbus Roman}
\fi
% Use upquote if available, for straight quotes in verbatim environments
\IfFileExists{upquote.sty}{\usepackage{upquote}}{}
\IfFileExists{microtype.sty}{% use microtype if available
  \usepackage[]{microtype}
  \UseMicrotypeSet[protrusion]{basicmath} % disable protrusion for tt fonts
}{}
\makeatletter
\@ifundefined{KOMAClassName}{% if non-KOMA class
  \IfFileExists{parskip.sty}{%
    \usepackage{parskip}
  }{% else
    \setlength{\parindent}{0pt}
    \setlength{\parskip}{6pt plus 2pt minus 1pt}}
}{% if KOMA class
  \KOMAoptions{parskip=half}}
\makeatother
\usepackage{xcolor}
\IfFileExists{xurl.sty}{\usepackage{xurl}}{} % add URL line breaks if available
\IfFileExists{bookmark.sty}{\usepackage{bookmark}}{\usepackage{hyperref}}
\hypersetup{
  pdftitle={Medial Septal Projections to the Parasubiculum},
  pdfauthor={Daniel Parthier},
  hidelinks,
  pdfcreator={LaTeX via pandoc}}
\urlstyle{same} % disable monospaced font for URLs
\usepackage[left=3cm, right=3cm, top=2.5cm, bottom=2.5cm]{geometry}
\usepackage{color}
\usepackage{fancyvrb}
\newcommand{\VerbBar}{|}
\newcommand{\VERB}{\Verb[commandchars=\\\{\}]}
\DefineVerbatimEnvironment{Highlighting}{Verbatim}{commandchars=\\\{\}}
% Add ',fontsize=\small' for more characters per line
\usepackage{framed}
\definecolor{shadecolor}{RGB}{248,248,248}
\newenvironment{Shaded}{\begin{snugshade}}{\end{snugshade}}
\newcommand{\AlertTok}[1]{\textcolor[rgb]{0.94,0.16,0.16}{#1}}
\newcommand{\AnnotationTok}[1]{\textcolor[rgb]{0.56,0.35,0.01}{\textbf{\textit{#1}}}}
\newcommand{\AttributeTok}[1]{\textcolor[rgb]{0.77,0.63,0.00}{#1}}
\newcommand{\BaseNTok}[1]{\textcolor[rgb]{0.00,0.00,0.81}{#1}}
\newcommand{\BuiltInTok}[1]{#1}
\newcommand{\CharTok}[1]{\textcolor[rgb]{0.31,0.60,0.02}{#1}}
\newcommand{\CommentTok}[1]{\textcolor[rgb]{0.56,0.35,0.01}{\textit{#1}}}
\newcommand{\CommentVarTok}[1]{\textcolor[rgb]{0.56,0.35,0.01}{\textbf{\textit{#1}}}}
\newcommand{\ConstantTok}[1]{\textcolor[rgb]{0.00,0.00,0.00}{#1}}
\newcommand{\ControlFlowTok}[1]{\textcolor[rgb]{0.13,0.29,0.53}{\textbf{#1}}}
\newcommand{\DataTypeTok}[1]{\textcolor[rgb]{0.13,0.29,0.53}{#1}}
\newcommand{\DecValTok}[1]{\textcolor[rgb]{0.00,0.00,0.81}{#1}}
\newcommand{\DocumentationTok}[1]{\textcolor[rgb]{0.56,0.35,0.01}{\textbf{\textit{#1}}}}
\newcommand{\ErrorTok}[1]{\textcolor[rgb]{0.64,0.00,0.00}{\textbf{#1}}}
\newcommand{\ExtensionTok}[1]{#1}
\newcommand{\FloatTok}[1]{\textcolor[rgb]{0.00,0.00,0.81}{#1}}
\newcommand{\FunctionTok}[1]{\textcolor[rgb]{0.00,0.00,0.00}{#1}}
\newcommand{\ImportTok}[1]{#1}
\newcommand{\InformationTok}[1]{\textcolor[rgb]{0.56,0.35,0.01}{\textbf{\textit{#1}}}}
\newcommand{\KeywordTok}[1]{\textcolor[rgb]{0.13,0.29,0.53}{\textbf{#1}}}
\newcommand{\NormalTok}[1]{#1}
\newcommand{\OperatorTok}[1]{\textcolor[rgb]{0.81,0.36,0.00}{\textbf{#1}}}
\newcommand{\OtherTok}[1]{\textcolor[rgb]{0.56,0.35,0.01}{#1}}
\newcommand{\PreprocessorTok}[1]{\textcolor[rgb]{0.56,0.35,0.01}{\textit{#1}}}
\newcommand{\RegionMarkerTok}[1]{#1}
\newcommand{\SpecialCharTok}[1]{\textcolor[rgb]{0.00,0.00,0.00}{#1}}
\newcommand{\SpecialStringTok}[1]{\textcolor[rgb]{0.31,0.60,0.02}{#1}}
\newcommand{\StringTok}[1]{\textcolor[rgb]{0.31,0.60,0.02}{#1}}
\newcommand{\VariableTok}[1]{\textcolor[rgb]{0.00,0.00,0.00}{#1}}
\newcommand{\VerbatimStringTok}[1]{\textcolor[rgb]{0.31,0.60,0.02}{#1}}
\newcommand{\WarningTok}[1]{\textcolor[rgb]{0.56,0.35,0.01}{\textbf{\textit{#1}}}}
\usepackage{longtable,booktabs}
% Correct order of tables after \paragraph or \subparagraph
\usepackage{etoolbox}
\makeatletter
\patchcmd\longtable{\par}{\if@noskipsec\mbox{}\fi\par}{}{}
\makeatother
% Allow footnotes in longtable head/foot
\IfFileExists{footnotehyper.sty}{\usepackage{footnotehyper}}{\usepackage{footnote}}
\makesavenoteenv{longtable}
\usepackage{graphicx,grffile}
\makeatletter
\def\maxwidth{\ifdim\Gin@nat@width>\linewidth\linewidth\else\Gin@nat@width\fi}
\def\maxheight{\ifdim\Gin@nat@height>\textheight\textheight\else\Gin@nat@height\fi}
\makeatother
% Scale images if necessary, so that they will not overflow the page
% margins by default, and it is still possible to overwrite the defaults
% using explicit options in \includegraphics[width, height, ...]{}
\setkeys{Gin}{width=\maxwidth,height=\maxheight,keepaspectratio}
% Set default figure placement to htbp
\makeatletter
\def\fps@figure{htbp}
\makeatother
\setlength{\emergencystretch}{3em} % prevent overfull lines
\providecommand{\tightlist}{%
  \setlength{\itemsep}{0pt}\setlength{\parskip}{0pt}}
\setcounter{secnumdepth}{5}
\usepackage{booktabs}
\usepackage{float}
\usepackage[format=plain,
labelfont={bf,small,it},
textfont={small,it}]{caption}
\renewcommand\chaptername{}
\pagestyle{plain}
\raggedbottom
\usepackage[]{natbib}
\bibliographystyle{apalike}

\title{Medial Septal Projections to the Parasubiculum}
\usepackage{etoolbox}
\makeatletter
\providecommand{\subtitle}[1]{% add subtitle to \maketitle
  \apptocmd{\@title}{\par {\large #1 \par}}{}{}
}
\makeatother
\subtitle{Inaugural-Dissertation\\
to obtain the academic degree\\
Doctor rerum naturalium (Dr.~rer. nat.)\\
~\\
submitted to the Department of Biology, Chemistry and Pharmacy\\
of Freie Universität Berlin by}
\author{Daniel Parthier}
\date{19.02.2021}

\begin{document}
\maketitle

{
\setcounter{tocdepth}{1}
\tableofcontents
}
\setstretch{1.5}
\hypertarget{prerequisites}{%
\chapter{Prerequisites}\label{prerequisites}}

The experimental work of this thesis was completed from xx.xx.xx to xx.xx.xx under the supervision of Prof.~Dr.~Dietmar Schmitz at the Neuroscience Research Centre (NWFZ) of the Charité.

\hypertarget{intro}{%
\chapter{Introduction}\label{intro}}

\hypertarget{theta}{%
\section{Theta}\label{theta}}

\hypertarget{medial-septum}{%
\section{Medial Septum}\label{medial-septum}}

The medial septum (MS) a structure located in the forebrain (Figure \ref{fig:MS-Scheme}) can be found across species.

\begin{figure}[H]

{\centering \includegraphics{Figures/Schemes/MS_schematic} 

}

\caption{Schematic of medial septum in saggital (A) and coronal (B) section. Blue shows the location of the MS, light gray represents ventricular space and dark gray fibre bundels. The dashed line in (A) indicates the position of the section of (B). Modified from Allen Brain Atlas.}\label{fig:MS-Scheme}
\end{figure}

\hypertarget{parasubiculum}{%
\section{Parasubiculum}\label{parasubiculum}}

The parasubiculum (PaS) on the other hand is part of the parahippocampal formation (Figure \ref{fig:PaS-Scheme}), the posterior part of the mouse brain, and is found in different species \citep{ding_comparative_2013}.
It is adjacent to the medial entorhinal cortex and presubiculum which, as the PaS, have spatially linked functional cell types coding for different aspects used for navigation (Citation needed).

\begin{figure}[H]
\includegraphics[width=0.9\linewidth,height=0.9\textheight]{Figures/Schemes/PaS_schematic} \caption{Schematic of parasubiculum in a saggital (A) and horizontal (B) section. turqoise shows the location of the PaS, light gray represents ventricular space and dark gray fibre bundels. The dark green in (B) represents the pyramidal cell layer II of the medial entorhinal cortex. The cornu ammonis and the dentate gyrus are marked in light blue.}\label{fig:PaS-Scheme}
\end{figure}

\hypertarget{spatial-navigation}{%
\section{Spatial Navigation}\label{spatial-navigation}}

\hypertarget{theta-1}{%
\subsection{Theta}\label{theta-1}}

You can label chapter and section titles using \texttt{\{\#label\}} after them, e.g., we can reference Chapter \ref{intro}. If you do not manually label them, there will be automatic labels anyway, e.g., Chapter \ref{methods}.

Figures and tables with captions will be placed in \texttt{figure} and \texttt{table} environments, respectively.

\begin{Shaded}
\begin{Highlighting}[]
\KeywordTok{par}\NormalTok{(}\DataTypeTok{mar =} \KeywordTok{c}\NormalTok{(}\DecValTok{4}\NormalTok{, }\DecValTok{4}\NormalTok{, }\FloatTok{.1}\NormalTok{, }\FloatTok{.1}\NormalTok{))}
\KeywordTok{plot}\NormalTok{(pressure, }\DataTypeTok{type =} \StringTok{'b'}\NormalTok{, }\DataTypeTok{pch =} \DecValTok{19}\NormalTok{)}
\end{Highlighting}
\end{Shaded}

\begin{figure}

{\centering \includegraphics[width=1\linewidth]{Thesis_files/figure-latex/nice-fig-1} 

}

\caption{Here is a nice figure!}\label{fig:nice-fig}
\end{figure}

Reference a figure by its code chunk label with the \texttt{fig:} prefix, e.g., see Figure \ref{fig:nice-fig}. Similarly, you can reference tables generated from \texttt{knitr::kable()}, e.g., see Table \ref{tab:nice-tab}.

\begin{Shaded}
\begin{Highlighting}[]
\NormalTok{knitr}\OperatorTok{::}\KeywordTok{kable}\NormalTok{(}
  \KeywordTok{head}\NormalTok{(iris, }\DecValTok{20}\NormalTok{), }\DataTypeTok{caption =} \StringTok{'Here is a nice table!'}\NormalTok{,}
  \DataTypeTok{booktabs =} \OtherTok{TRUE}
\NormalTok{)}
\end{Highlighting}
\end{Shaded}

\begin{table}

\caption{\label{tab:nice-tab}Here is a nice table!}
\centering
\begin{tabular}[t]{rrrrl}
\toprule
Sepal.Length & Sepal.Width & Petal.Length & Petal.Width & Species\\
\midrule
5.1 & 3.5 & 1.4 & 0.2 & setosa\\
4.9 & 3.0 & 1.4 & 0.2 & setosa\\
4.7 & 3.2 & 1.3 & 0.2 & setosa\\
4.6 & 3.1 & 1.5 & 0.2 & setosa\\
5.0 & 3.6 & 1.4 & 0.2 & setosa\\
\addlinespace
5.4 & 3.9 & 1.7 & 0.4 & setosa\\
4.6 & 3.4 & 1.4 & 0.3 & setosa\\
5.0 & 3.4 & 1.5 & 0.2 & setosa\\
4.4 & 2.9 & 1.4 & 0.2 & setosa\\
4.9 & 3.1 & 1.5 & 0.1 & setosa\\
\addlinespace
5.4 & 3.7 & 1.5 & 0.2 & setosa\\
4.8 & 3.4 & 1.6 & 0.2 & setosa\\
4.8 & 3.0 & 1.4 & 0.1 & setosa\\
4.3 & 3.0 & 1.1 & 0.1 & setosa\\
5.8 & 4.0 & 1.2 & 0.2 & setosa\\
\addlinespace
5.7 & 4.4 & 1.5 & 0.4 & setosa\\
5.4 & 3.9 & 1.3 & 0.4 & setosa\\
5.1 & 3.5 & 1.4 & 0.3 & setosa\\
5.7 & 3.8 & 1.7 & 0.3 & setosa\\
5.1 & 3.8 & 1.5 & 0.3 & setosa\\
\bottomrule
\end{tabular}
\end{table}

look at figure \ref{fig:nice-fig}

You can write citations, too. For example, we are using the \textbf{bookdown} package \citep{R-bookdown} in this sample book, which was built on top of R Markdown and \textbf{knitr} \citep{xie2015}.

\hypertarget{materials-and-methods}{%
\chapter{Materials and Methods}\label{materials-and-methods}}

\hypertarget{in-vitro}{%
\section{In-Vitro}\label{in-vitro}}

\hypertarget{slice-preparation}{%
\subsection{Slice Preparation}\label{slice-preparation}}

For slices animals (n = XX PV-Cre mice, n = XX ChAT-Cre mice ) were deeply anaesthetised using isoflurane decapitated, the brain removed and quickly transferred to ice-cold slicing solution. sucrose artificial cerebral spinal fluid (SACF). SACF contained 87 NaCl, 26 NaHCO\textsubscript{3}, 10 Glucose, 50 Sucrose, 2.5 KCl, 1.25 NaH\textsubscript{2}PO\textsubscript{4}, 0.5 CaCl\textsubscript{2}, 3 MgCl\textsubscript{2} · 6H\textsubscript{2}O.
400 µm horizontal slices were produced using a vibratome (VT1200S, Leica Biosystems, Wetzlar, Germany), then transferred and stored in an interface chamber for up to 1-6 h. Slices were perfused with ACSF (119 NaCl, 26 NaHCO\textsubscript{3}, 10 Glucose, 2.5 KCl, 1 NaH\textsubscript{2}PO\textsubscript{4}, 2.5 CaCl\textsubscript{2}, 1.3 MgCl\textsubscript{2} · 6H\textsubscript{2}O) and oxygenated during the whole period.

\hypertarget{slice-recordings}{%
\subsection{Slice Recordings}\label{slice-recordings}}

Cells were identified using XXXX DIC XXXX and recorded in the whole-cell patch-clamp configuration using a glass electrode filled with intracellular solution (120 K-Gluconate, 10 Hepes, 10 KCl, 5 EGTA, 2 MgSO\textsubscript{4} · 7H\textsubscript{2}, 3 MgATP, 1 NaGTP, 5 Phosphocreatine Na, 0.2\% Biocytin) to record currents, voltage and later identify cells via Biocytin staining. After opening, cells were hold in voltage-clamp at -60 mV to measure series resistance. Then cells were switched to current clamp to perform a characterisation of cell properties (current steps of 40-100 pA). To assess septal connectivity to the parasubiculum Channelrhodopsin expressing fibres in the area were activated using 10 ms light pulses. For the light stimulation a range of the pulse frequencies (10, 20, 40 Hz) was used to evoke membrane potential changes (EPSPs) or membrane current (EPSCs). To detect masked inputs due to low driving force or simultaneously occuring synaptic inputs the holding potential of the cell was changed from -60, to -80 and then to -50 mV after at least 10 trials.

all in c/mM

Drugs used:

\hypertarget{results}{%
\chapter{Results}\label{results}}

Here is a review of existing methods.

\hypertarget{discussion}{%
\chapter{Discussion}\label{discussion}}

Here is a review of existing methods.

  \bibliography{book.bib,packages.bib,references.bib}

\end{document}
